\begin{textsample}{NEG Dim 5 – human – Score -10.00 – rj/rj\_ef\_matematica\_2\_p1.txt}  \label{ex:f5_neg_012}
No Ensino Fundamental, a área de Matemática está presente articulando seus eixos temáticos.
Desta forma a BNCC propõe o ensino da Matemática dividido em Unidades Temáticas, com características de semelhança para um ensinamento de progressão contínua.
Para o \textbf{desenvolvimento} das \textbf{habilidades} prevista na \textbf{Base} \textbf{Nacional} Comum \textbf{Curricular}, é imprescindível levar em consideração os \textbf{conhecimentos} matemáticos e as experiências do \textbf{estudante}.
A aprendizagem está \textbf{relacionada} à compreensão de sentido dos \textbf{objetos} matemáticos.
Esses sentidos são consequências das conexões que os \textbf{estudantes} formam entre os \textbf{objetos} e seu cotidiano, entre eles e os diferentes temas matemáticos e, por fim, entre eles e os demais componentes \textbf{curriculares}.
A Matemática é uma peça importante na construção da cidadania, nos \textbf{conhecimentos} científicos e recursos tecnológicos.
O seu ensino deve ser de construção do pensamento lógico-matemático, despertando o espírito investigativo, procurando desenvolver nos alunos competências para compreender e transformar a realidade, de maneira prazerosa.

% matched lemmas: base, conhecimento, curricular, desenvolvimento, estudante, habilidade, nacional, objeto, relacionar
\end{textsample}
